%% toc: P3 36 36.8,36.9
\chapter{Definable Second-Order Quantifiers and the Comparison of Theories}
\label{sec:p03-w03-p03-046-dsoq}

A first-order theory \(T\) can be strengthened by allowing quantification
over relations of a prescribed kind on its models: over all subsets
(monadic logic), over all permutations of order at most two
(permutational logic), or over all binary relations (second-order logic).
Comparing such strengthened theories by interpretation gives a measure of
how complicated the models of \(T\) are. This chapter presents the
definitions on which that comparison rests, a group of sufficient
conditions under which full second-order logic of an infinite set can be
interpreted in the monadic or the permutational theory of \(T\), and the
basic theory of definable second-order quantifiers over pure sets: their
mutual interpretations, the separation of the first three of the four
principal quantifiers, and the simple-definability criteria that drive the
classification.

\emph{Source and version.} Two works are followed, each with its own
locators. Sections~\ref{sec:p03-w03-p03-046-quantifiers}
to~\ref{sec:p03-w03-p03-046-sufficient} follow the paper of Baldwin and
Shelah on second-order quantifiers and the complexity of theories, in the
published version of 1985 as made available on Shelah's archive
\cite{BaldwinShelah1985}, Sections~1.3, 1.4, 2.1 and 2.2 (pp.~234--244).
\Cref{sec:p03-w03-p03-046-baldwin} follows Section~1 of Baldwin's chapter
\emph{Definable Second-Order Quantifiers}, Chapter~XII of
\emph{Model-Theoretic Logics}, in the 1985 Springer-Verlag printing
(pp.~445--477) as reproduced open access on Project Euclid
\cite{Baldwin1985DSOQ}; the 2017 Cambridge reissue of the collection
\cite{BarwiseFeferman1985} was not the version read, and every chapter
locator below is a page of the 1985 printing.

\emph{Inputs.} A consistent first-order theory \(T\) in a language \(L\)
with no finite models; a quantifier condition \(\varphi(\bar r)\) in pure
identity theory; formulas of \(L\) that code pairs.
\emph{Conclusions.} Precise notions of interpretation between theories in
logics with second-order quantifiers; monotonicity of Hanf and
L\"owenheim numbers under spectrum-preserving interpretations; three
sufficient conditions (a definable pairing function, a coding of a
product, a definable equivalence relation with infinitely many infinite
classes) for interpreting second-order logic; and the interpretation of
second-order logic in the permutational theory of linear order.

\emph{Metatheory.} ZFC. Choice is used through first-order compactness
(\cref{thm:p03-compactness-ultraproduct}) and in partitioning an infinite set into
two-element blocks.

\section{Second-order quantifiers and the four principal logics}
\label{sec:p03-w03-p03-046-quantifiers}

\begin{hbdefinition}{def:p03-w03-p03-046-quantifier}{Second-order quantifier \(Q_{\varphi(\bar r)}\)}
Let \(\varphi(\bar r)\) be a first-order sentence whose only nonlogical
symbols are relation variables \(\bar r=\langle r_0,\dots,r_{n-1}\rangle\)
and equality. For a language \(L\), the logic \(L(Q_{\varphi(\bar r)})\)
extends first-order logic by the following clause. Let \(M\) be an
\(L\)-structure and \(L'=L(\bar r)\). For a formula \(\psi(\bar r)\) of
\(L'(Q_{\varphi(\bar r)})\),
\(M\entails Q_{\varphi(\bar r)}\,\psi(\bar r)\) holds iff there are
relations \(R_0,\dots,R_{n-1}\) on \(M\), of the arities of \(\bar r\),
with \(M\entails\varphi(\bar R)\wedge\psi(\bar R)\).
For a first-order theory \(T\), the pair \((T,Q_{\varphi(\bar r)})\) is the
set of \(L(Q_{\varphi(\bar r)})\)-sentences true in every model of \(T\)
\cite[p.~239, Definition~2.1.1]{BaldwinShelah1985}.
\end{hbdefinition}

All relation variables of one quantifier are taken to have a common arity
\(k\), as in the source, which remarks that formulas with mixed arities can
be reduced to this case \cite[p.~239]{BaldwinShelah1985}; that reduction
is not used below.

\begin{hbdefinition}{def:p03-w03-p03-046-four}{\(R_\varphi(A)\) and the four logics}
Let \(\varphi(\bar s)=\varphi(s_0,\dots,s_{n-1})\) be a formula of pure
identity theory in which each \(s_i\) is a \(k\)-ary relation variable. For
a set \(A\), \(R_\varphi(A)\) is the set of tuples
\(\langle S_0,\dots,S_{n-1}\rangle\) of \(k\)-ary relations on \(A\) with
\(A\entails\varphi(S_0,\dots,S_{n-1})\); one writes \(\bar S\in
R_\varphi(A)\) for \(A\entails\varphi(\bar S)\)
\cite[p.~240, Notation~2.1.2]{BaldwinShelah1985}. The four principal
cases are:
\begin{enumerate}
  \item \((T,\mathrm{1st})\), the first-order theory of the models of \(T\);
  \item \((T,\mathit{Mon})\), the monadic theory: quantification over all
        unary relations (subsets);
  \item \((T,\text{1-1})\), the permutational theory:
        \(M\entails Q f\,\varphi(f)\) holds exactly when there is a
        permutation \(\alpha\) of \(M\) of order two with
        \(M\entails\varphi(\alpha)\), that is, a bijection with
        \(\alpha\circ\alpha=\mathrm{id}_M\) (an \emph{involution});
        the relation quantified over is the graph of \(\alpha\);
  \item \((T,\mathrm{2nd})\), the second-order theory: quantification over
        all binary relations.
\end{enumerate}
\(T_\infty\) denotes the theory, in the language with equality only, of an
infinite set, and \(\mathit{Th}(<)\) the first-order theory of infinite
linear orders \cite[pp.~231, 240]{BaldwinShelah1985}.
\end{hbdefinition}

The source explains the permutational case as quantification over
equivalence relations each of whose classes has at most two elements, the
classes being the orbits of the involution \cite[p.~240]{BaldwinShelah1985}.
For an involution \(\alpha\) write
\(\operatorname{Mov}(\alpha)=\{x:\alpha(x)\neq x\}\); this set is
first-order definable from the graph of \(\alpha\). Every set whose
cardinality is infinite is \(\operatorname{Mov}(\alpha)\) for some
involution \(\alpha\) of any structure containing it: an infinite set
\(X\) has \(|X|=|X\times 2|\), so it splits into two-element blocks, and
\(\alpha\) exchanges the members of each block and fixes everything outside
\(X\).

\section{Interpretations, spectra and Hanf numbers}
\label{sec:p03-w03-p03-046-interpretations}

\begin{hbdefinition}{def:p03-w03-p03-046-spectrum}{Spectrum, Hanf and L\"owenheim numbers relative to \(T\)}
For a sentence \(\phi\) of a logic \(\mathscr L\),
\(\operatorname{spec}(\phi)=\{\kappa\ge\aleph_0:\exists M\,(M\entails\phi,\
|M|=\kappa)\}\), and for a theory \(T\),
\(\operatorname{Spec}_T(\phi)=\{\kappa\ge\aleph_0:\exists M\,(M\entails
T\cup\{\phi\},\ |M|=\kappa)\}\); \(\phi\) is \emph{bounded} if its spectrum
is a set \cite[p.~234, Notation~1.3.1]{BaldwinShelah1985}. Finite models
are ignored throughout. The \emph{Hanf number of \(\mathscr L\) relative to
\(T\)}, \(h^T_{\mathscr L}\), is the least \(\kappa\) such that for every
\(\mathscr L\)-sentence \(\phi\) in the language of \(T\), if
\(T\cup\{\phi\}\) has a model of power \(\kappa\) then it has arbitrarily
large models. The \emph{L\"owenheim number of \(\mathscr L\) relative to
\(T\)}, \(ls^T_{\mathscr L}\), is the least \(\kappa\) such that for every
\(\mathscr L\)-sentence \(\phi\), if \(T\cup\{\phi\}\) has a model then it
has a model of cardinality at most \(\kappa\)
\cite[p.~235, Definition~1.3.2]{BaldwinShelah1985}.
\end{hbdefinition}

\begin{hbdefinition}{def:p03-w03-p03-046-syntactic}{Syntactic and spectrum-preserving interpretations}
Let \(T_1\) be an \(\mathscr L_1\)-theory and \(T_2\) an
\(\mathscr L_2\)-theory.
\begin{enumerate}
  \item \(T_1\) is \emph{syntactically interpretable} in \(T_2\) if there
        is a map \(*\) assigning to each \(\mathscr L_1\)-sentence \(\phi\)
        an \(\mathscr L_2\)-sentence \(\phi^*\) with \(T_1\vdash\phi\) iff
        \(T_2\vdash\phi^*\) \cite[p.~237, Definition~1.4.1]{BaldwinShelah1985}.
  \item A \emph{spectrum-preserving interpretation} of \(T_1\) into \(T_2\)
        is a map \(*\) from \(\mathscr L_1\)-sentences to
        \(\mathscr L_2\)-sentences with
        \(\operatorname{Spec}_{T_1}(\phi)=\operatorname{Spec}_{T_2}(\phi^*)\)
        for every \(\mathscr L_1\)-sentence \(\phi\)
        \cite[p.~237, Definition~1.4.2]{BaldwinShelah1985}.
\end{enumerate}
\end{hbdefinition}

The first notion places no control on the cardinalities of models, and the
source notes that it has no effect on Hanf or L\"owenheim numbers
\cite[p.~237]{BaldwinShelah1985}. The second notion does control them.

\begin{hblemma}{thm:p03-w03-p03-046-spectrum-monotone}{proved=annex:p03-w03-p03-046-spectrum-monotone}{Spectrum-preserving interpretations}
If there is a spectrum-preserving interpretation of the
\(\mathscr L_1\)-theory \(T_1\) into the \(\mathscr L_2\)-theory \(T_2\),
then \(ls^{T_1}_{\mathscr L_1}\le ls^{T_2}_{\mathscr L_2}\) and
\(h^{T_1}_{\mathscr L_1}\le h^{T_2}_{\mathscr L_2}\), whenever the numbers
on the right exist \cite[p.~237, Lemma~1.4.3]{BaldwinShelah1985}.
\end{hblemma}

The notion of interpretation used for the classification is stronger: it
interprets the quantifier, uniformly over all models of \(T\) and all
infinite cardinalities.

\begin{hbdefinition}{def:p03-w03-p03-046-interpretation}{First-order interpretation of a second-order quantifier}
Let \(\varphi(r_0,\dots,r_{m-1})\) be a formula in \(m\) \(k\)-ary relation
variables, \(\psi(\bar s)\) a formula in \(m'\) \(k'\)-ary relation
variables, and \(T\) a first-order theory with no finite models.
\((T_\infty,Q_{\varphi(\bar r)})\) is \emph{first-order interpretable} in
\((T,Q_{\psi(\bar s)})\), written
\((T_\infty,Q_{\varphi(\bar r)})\le(T,Q_{\psi(\bar s)})\), if there are
\(n<\omega\) and first-order formulas
\begin{gather*}
\pi(\bar y,\bar s_0,\dots,\bar s_{n-1}),\qquad
\theta(x,\bar y,\bar s_0,\dots,\bar s_{n-1}),\\
\chi_i(x_0,\dots,x_{k-1},\bar y,\bar s_0,\dots,\bar s_{n-1},\\
\bar y',\bar s'_0,\dots,\bar s'_{n-1})\quad(i<m),
\end{gather*}
each \(\bar s_j,\bar s'_j\) an \(m'\)-tuple of \(k'\)-ary relation
variables, such that:
\begin{enumerate}
  \item if \(M\entails T\), \(\bar a\in M\), \(\bar S_0,\dots,\bar
        S_{n-1}\in R_\psi(M)\) and \(M\entails\pi(\bar a,\bar S_0,\dots,\bar
        S_{n-1})\), then \(B=\theta(M,\bar a,\bar S_0,\dots,\bar S_{n-1})\) is
        an infinite set and \(R_\varphi(B)\) is exactly the set of tuples
        \[\langle\chi_0(M,\bar a,\bar S,\bar a',\bar S'),\dots,
        \chi_{m-1}(M,\bar a,\bar S,\bar a',\bar S')\rangle\] with
        \(\bar a'\in M\) and \(\bar S'_0,\dots,\bar S'_{n-1}\in R_\psi(M)\);
  \item for every \(\kappa\ge\aleph_0\) there are \(M\entails T\),
        \(\bar a\in M\) and \(\bar S_0,\dots,\bar S_{n-1}\in R_\psi(M)\)
        with \(M\entails\pi(\bar a,\bar S_0,\dots,\bar S_{n-1})\) and
        \(|\theta(M,\bar a,\bar S_0,\dots,\bar S_{n-1})|=\kappa\).
\end{enumerate}
If each of two such pairs is interpretable in the other they are
\emph{bi-interpretable}, written \(\equiv\)
\cite[p.~241, Definition~2.1.3(a), (c)]{BaldwinShelah1985}.
\end{hbdefinition}

\begin{hbconvention}{def:p03-w03-p03-046-reading}{Reading of the coding parameters}
In clause~(1) the source writes the coding formulas with the primed
parameters \(\bar a',\bar S'\) ranging over all of \(M\) and
\(R_\psi(M)\) \cite[p.~241]{BaldwinShelah1985}. Here the formulas
\(\chi_i\) carry the fixed parameters \(\bar a,\bar S\) that define \(B\)
as well as the varying coding parameters \(\bar a',\bar S'\); a formula
that ignores one block is a special case. Every proof below exhibits
\(\pi\), \(\theta\) and \(\chi_i\) explicitly and verifies both inclusions
of clause~(1). The source also records that one inclusion suffices after a
modification of \(\pi\) \cite[p.~242]{BaldwinShelah1985}; that remark is
not used.
\end{hbconvention}

For the target \(\mathrm{2nd}\) the quantifier condition is trivially true
of one binary relation variable, so \(R_{\mathrm{2nd}}(B)\) is the set of
all binary relations on \(B\). Clause~(1) then says that, as the coding
parameters \(\bar a',\bar S'\) vary, the sets
\(\chi(M,\bar a,\bar S,\bar a',\bar S')\) are exactly the binary relations on \(B\). For \(\mathit{Mon}\) the parameters
are arbitrary subsets of \(M\); for \(\text{1-1}\) they are graphs of
involutions of \(M\).

The source extends the notion to an interpretation of \((T_0,Q_\varphi)\)
in \((T_1,Q_\psi)\) that also carries the relations of a first-order theory
\(T_0\), and to classes of structures and arbitrary logics, written
\((K_0,\mathscr L_1)\le_{\mathscr L}(K_1,\mathscr L_2)\) when the
interpreting formulas lie in \(\mathscr L\)
\cite[p.~241, Definition~2.1.3(b), (d)]{BaldwinShelah1985}; it calls such
an interpretation \emph{almost decreasing} if for each model \(M\) the
cardinal \(|M|\) is bounded in terms of \(|\theta(M,\bar a,\bar S)|\)
\cite[p.~242, Definition~2.1.4]{BaldwinShelah1985}. Only clause~(a) is used
in this chapter.

\begin{proofboundary}{def:p03-w03-p03-046-only-four}{The four-quantifier classification}
Baldwin and Shelah report as a theorem of Shelah that
\((T_\infty,\mathrm{1st})<(T_\infty,\mathit{Mon})<(T_\infty,\text{1-1})<(T_\infty,\mathrm{2nd})\),
the inequalities being strict, and that for every quantifier
\(Q_{\psi(\bar r)}\) the pair \((T_\infty,Q_{\psi(\bar r)})\) is
first-order bi-interpretable with one of these four
\cite[p.~242, Theorem~2.1.6]{BaldwinShelah1985}. No proof of this
classification is given in this chapter, and none of the results below
depends on it. The classification explains why the four logics of
\cref{def:p03-w03-p03-046-four} are the ones compared.
\end{proofboundary}

Since \((T_\infty,\mathrm{2nd})\not\le(T_\infty,\mathit{Mon})\) by the
classification, an interpretation \((T_\infty,\mathrm{2nd})\le(T,\mathit{Mon})\)
records complexity contributed by \(T\) itself. Baldwin and Shelah take
such an interpretation as the mark of a theory beyond analysis and aim to
show that \((T,\mathit{Mon})\) is simple when none exists
\cite[p.~242]{BaldwinShelah1985}. The next section gives the sufficient
conditions for interpretations from the source's Section~2.2.

\section{Sufficient conditions for interpreting second-order logic}
\label{sec:p03-w03-p03-046-sufficient}

\begin{hbdefinition}{def:p03-w03-p03-046-pairing}{Definable pairing function; coding}
\begin{enumerate}
  \item \(T\) has a \emph{definable pairing function} if there are
        formulas \(\rho(x,y,z)\), \(\rho_0(x,y)\), \(\rho_1(x,y)\) which, in
        every model of \(T\), define the graphs of functions \(\rho(x,y)\),
        \(\rho_0(x)\), \(\rho_1(x)\) with
        \(\rho(\rho_0(x),\rho_1(x))=x\), \(\rho_0(\rho(x,y))=x\) and
        \(\rho_1(\rho(x,y))=y\)
        \cite[p.~242, Definition~2.2.1]{BaldwinShelah1985}.
  \item \(T\) \emph{admits coding} if there is a model \(M\) of \(T\)
        containing infinite sets \(B_0,B_1,C\) such that some first-order
        formula \(\phi(x_0,x_1,y,\bar z,B_0,B_1,C)\), in which
        \(B_0,B_1,C\) occur as unary predicates, defines, for suitable
        constants \(\bar d\) substituted for \(\bar z\), a one-to-one
        function from \(B_0\times B_1\) onto \(C\).
  \item \(T\) \emph{admits strong coding} if there is such a formula
        \(\phi(x_1,x_2,\bar y,\bar z)\) without unary predicates, defining for
        some \(M\), definable infinite \(B_0,B_1,C\) and parameters \(\bar d\)
        a one-to-one function from \(B_0\times B_1\) onto \(C\)
        \cite[pp.~242--243, Definition~2.2.3]{BaldwinShelah1985}.
\end{enumerate}
\end{hbdefinition}

\begin{hblemma}{thm:p03-w03-p03-046-pairing}{proved=annex:p03-w03-p03-046-pairing}{A definable pairing function interprets second-order logic monadically}
If \(T\) is a consistent first-order theory with no finite models and
\(T\) has a definable pairing function, then
\((T_\infty,\mathrm{2nd})\le(T,\mathit{Mon})\)
\cite[p.~242, Lemma~2.2.2]{BaldwinShelah1985}.
\end{hblemma}

\begin{hblemma}{thm:p03-w03-p03-046-coding}{proved=annex:p03-w03-p03-046-coding}{Coding and monadic logic}
If \(T\) is a first-order theory with no finite models and \(T\) admits
coding, then \((T_\infty,\mathrm{2nd})\le(T,\mathit{Mon})\)
\cite[p.~243, Lemma~2.2.4]{BaldwinShelah1985}.
\end{hblemma}

The source observes that a theory may admit coding without having a
first-order definable pairing function, citing Lachlan for the fact that
no superstable theory has a first-order pairing function while some
superstable theories admit coding \cite[p.~243]{BaldwinShelah1985}; that
fact is not proved here. The hypothesis on \(\phi\) in
\cref{thm:p03-w03-p03-046-coding} is first-order, which is what makes the
interpretation uniform over all models.

For the permutational theory the parameters are involutions rather than
sets, and the relevant combinatorial hypothesis is an equivalence relation.

\begin{hblemma}{thm:p03-w03-p03-046-equiv-uniform}{proved=annex:p03-w03-p03-046-equiv-uniform}{Large equivalence relations and permutational logic}
Let \(T\) be a first-order theory with no finite models, and let
\(\phi(x,y,\bar z,s)\) be a first-order formula of \(L\) together with one
binary relation variable \(s\). Suppose that for every
\(\kappa\ge\aleph_0\) there are \(M\entails T\), \(\bar c\in M\), an
involution \(\sigma\) of \(M\) and a set \(A\subseteq M\) such that
\(\phi(x,y,\bar c,\sigma)\) defines on \(A\) an equivalence relation with at
least \(\kappa\) classes each of cardinality at least \(\kappa\). Then
\((T_\infty,\mathrm{2nd})\le(T,\text{1-1})\).
\end{hblemma}

\Cref{thm:p03-w03-p03-046-equiv-uniform} is the form in which the next two
results are proved here; it allows the defining formula a permutation
parameter, which the source uses in its proof for linear orders.

\begin{hblemma}{thm:p03-w03-p03-046-equiv}{proved=annex:p03-w03-p03-046-equiv}{An equivalence relation with infinitely many infinite classes}
If \(T\) is a first-order theory with no finite models and there is a
first-order formula \(\phi(x,y)\) such that for some \(M\entails T\) and some
\(A\subseteq M\), \(\phi(x,y)\) defines an equivalence relation on \(A\)
with infinitely many infinite classes, then
\((T_\infty,\mathrm{2nd})\le(T,\text{1-1})\)
\cite[p.~244, Lemma~2.2.6]{BaldwinShelah1985}.
\end{hblemma}

\begin{hbcorollary}{thm:p03-w03-p03-046-order}{proved=annex:p03-w03-p03-046-order}{Second-order logic in the permutational theory of order}
\((T_\infty,\mathrm{2nd})\le(\mathit{Th}(<),\text{1-1})\)
\cite[p.~244, Corollary~2.2.7]{BaldwinShelah1985}.
\end{hbcorollary}

The source proves \cref{thm:p03-w03-p03-046-equiv} by finding two subsets
of size \(\kappa\) coded by a third and appealing to its
Corollary~2.2.5, and sketches \cref{thm:p03-w03-p03-046-order} by defining
the blocks of a \(\kappa\)-dense order from a set of points
\cite[p.~244]{BaldwinShelah1985}. The annex proofs instead go through
\cref{thm:p03-w03-p03-046-equiv-uniform}, which builds the involutions
directly. Corollary~2.2.5 of the source, which states the coding criterion
for interpreting formulas in \(\mathit{Mon}\) or \(\text{1-1}\)
\cite[p.~243]{BaldwinShelah1985}, is not presented here.

\section{Definable second-order quantifiers over pure sets}
\label{sec:p03-w03-p03-046-baldwin}

This section follows Section~1 of Baldwin's chapter
\cite[pp.~445--451]{Baldwin1985DSOQ} in the 1985 printing. Baldwin's
notation differs from that of the preceding sections, and it is kept here:
the quantifier condition is written \(\psi\), the defining formulas
\(\phi\), \(\theta_0\), \(\theta\), \(\chi\), and the four principal
quantifiers \(Q_{\mathrm I}\), \(Q_{\mathrm{mon}}\), \(Q_{1-1}\),
\(Q_{\mathrm{II}}\). Baldwin's \(Q_{1-1}\) quantifies over equivalence
relations whose classes have two elements, while the logic
\(\text{1-1}\) of \cref{def:p03-w03-p03-046-four} quantifies over
involutions; \cref{thm:p03-w03-p03-046-b-permutations} shows that the two
are interchangeable.

\emph{Purpose.} The chapter asks which restrictions of second-order
quantification, cut out by a condition in pure identity theory, give
genuinely different logics, and compares them by first-order
interpretation \cite[pp.~445--446]{Baldwin1985DSOQ}.
\emph{Inputs.} A quantifier condition \(\psi(r)\) in pure identity theory;
infinite sets as models of the theory of equality.
\emph{Conclusions.} The interpretation preorder on definable quantifiers,
the chain \(Q_{\mathrm I}\le Q_{\mathrm{mon}}\le Q_{1-1}\le Q_{\mathrm{II}}\),
the strictness of its first two links, two simple-definability criteria
(\cref{thm:p03-w03-p03-046-b-mon,thm:p03-w03-p03-046-b-perm-crit}), and two
sufficient conditions for interpreting second-order logic in a theory
(\cref{thm:p03-w03-p03-046-b-codable,thm:p03-w03-p03-046-b-equiv}).

\subsection{Quantifiers, theories and conventions}
\label{sec:p03-w03-p03-046-b-conventions}

Baldwin works with theories rather than logics because the behaviour of a
logic depends heavily on the vocabulary: with only unary predicates,
monadic logic has Hanf and L\"owenheim numbers \(\aleph_0\)
(\cref{thm:p03-w03-p03-046-b-mso-sets}), while a binary function symbol
required to be a pairing function turns monadic logic into full
second-order logic (\cref{thm:p03-w03-p03-046-pairing})
\cite[p.~446]{Baldwin1985DSOQ}. Individual variables are \(x,y,z\),
relation variables \(r,s,t\), relations \(R,S,T\) and individuals
\(a,b,c\); \(\bar a\in A\) and \(\bar R\in A\) abbreviate tuples of
individuals and of relations on \(A\). For a formula \(\phi(\bar x,\bar
y,\bar s)\), \(\phi(A,\bar b,\bar S)=\{\bar a\in A:A\entails\phi(\bar
a,\bar b,\bar S)\}\). Equality is a logical symbol, and
\((\exists^{<k}x)\phi(x,\bar y,\bar r)\) abbreviates
\[
(\forall x_0)\cdots(\forall x_k)\Bigl(\bigwedge_{i<k+1}\phi(x_i,\bar y,\bar r)
\to\bigvee_{i<j<k+1}x_i=x_j\Bigr)
\]
\cite[p.~446]{Baldwin1985DSOQ}; it says there are at most \(k\) solutions.

\begin{hbdefinition}{def:p03-w03-p03-046-b-quantifier}{Definable second-order quantifier \(Q_{\psi(r)}\)}
Let \(\psi(r)\) be a formula whose only nonlogical symbol is the \(n\)-ary
relation variable \(r\). For each infinite set \(A\),
\(\mathscr R_\psi(A)\) is the collection of \(n\)-ary relations \(R\) on
\(A\) with \(A\entails\psi(R)\); the same notation is used when \(\psi\)
contains a finite sequence \(\bar r\) of relation variables
\cite[p.~447, Notation~1.2.1]{Baldwin1985DSOQ}. If \(\psi(r)\) is built
from \(r\), \(=\), first-order quantifiers and connectives only, then
\(Q_{\psi(r)}\) is the second-order quantifier with
\[
M\entails Q_{\psi(r)}\,\phi(r)\iff(\exists R)\in\mathscr R_\psi(M),\
M\entails\phi(R)
\]
\cite[p.~447, Definition~1.2.2]{Baldwin1985DSOQ}. For a first-order theory
\(T\), \((T,Q_\psi)\) is the collection of all \(Q_\psi\)-sentences of
\(L(T)\) valid on the models of \(T\)
\cite[p.~447, Definition~1.2.3]{Baldwin1985DSOQ}. By convention \(Q_\psi\)
alone means \((\mathrm{Th}(=),Q_\psi)\), \(\mathrm{Th}(=)\) being the theory of
equality, and \((\mathrm{Th}(<),Q_\psi)\) is the \(Q_\psi\)-theory of order
\cite[p.~447]{Baldwin1985DSOQ}.
\end{hbdefinition}

The first-order theory \(T_\psi\) with the single nonlogical symbol \(R\)
and the single axiom \(\psi(R)\) is attached to each quantifier, but it
carries less information than the quantifier, since a formula of the
extended logic may use several instances of the quantified relation
\cite[p.~447]{Baldwin1985DSOQ}. Models of \(\mathrm{Th}(=)\) are taken to be
infinite sets throughout, as in Notation~1.2.1.

\begin{hbdefinition}{def:p03-w03-p03-046-b-four}{The four principal quantifiers}
\begin{enumerate}
  \item \emph{First-order quantification} \(Q_{\mathrm I}\). Baldwin
        records that it is a definable second-order quantifier
        \cite[p.~447]{Baldwin1985DSOQ}; the condition used for it here is
        \(\psi_{\mathrm I}(r):=\exists x\forall y\,(r(y)\leftrightarrow
        y=x)\) with \(r\) unary, whose relations are the singletons, so
        that quantifying over them is quantifying over elements.
  \item \emph{Monadic quantification} \(Q_{\mathrm{mon}}\): \(r\) unary
        and \(\psi(r)\) any valid formula; \(Q_{\psi(r)}\) is then the
        monadic second-order quantifier \cite[p.~447]{Baldwin1985DSOQ}.
  \item \emph{Permutational quantification} \(Q_{1-1}\): \(r\) binary and
        \(\psi(r)\) says that \(r\) is an equivalence relation every class
        of which has two elements \cite[p.~447]{Baldwin1985DSOQ}. Its
        relations are called \emph{pairings} here.
  \item \emph{Full second-order quantification} \(Q_{\mathrm{II}}\),
        ranging over arbitrary relations \cite[pp.~445, 447]{Baldwin1985DSOQ}.
        Here it is taken with \(r\) binary and \(\psi\) valid, as in
        \cref{def:p03-w03-p03-046-four}; by
        \cref{thm:p03-w03-p03-046-b-maximal} this binary quantifier already
        interprets every definable quantifier, of any arity, whose
        condition holds of some relations on some infinite set.
\end{enumerate}
\end{hbdefinition}

Quantification over the automorphisms of an \(L\)-structure is not a
definable second-order quantifier, since being an automorphism cannot be
expressed without the symbols of \(L\) \cite[p.~447]{Baldwin1985DSOQ}.

\begin{hbdefinition}{def:p03-w03-p03-046-b-interpretation}{First-order interpretation of quantifiers}
Let \(r\) be \(k\)-ary. \(Q_{\psi(r)}\) is \emph{(first-order)
interpretable} in \((T,Q_{\phi(s)})\) if there are first-order formulas
\(\theta_0(x_0,\bar y,\bar s)\), \(\theta(x_0,\dots,x_{k-1},\bar y,\bar s)\)
and \(\chi(\bar x,\bar s)\) such that
\begin{enumerate}
  \item[(i)] if \(A\entails\chi(\bar a_0,\bar S_0)\), then
        \(\theta_0(A,\bar a_0,\bar S_0)\) is infinite,
        \(\bar S_0\in\mathscr R_\phi(A)\), and
        \[
        \bigl(\theta_0(A,\bar a_0,\bar S_0),\{\theta(A,\bar a,\bar S):\bar
        a\in A,\ \bar S\in\mathscr R_\phi(A)\}\bigr)
        =\bigl(\theta_0(A,\bar a_0,\bar S_0),\mathscr R_\psi(\theta_0(A,\bar
        a_0,\bar S_0))\bigr);
        \]
  \item[(ii)] for every infinite \(B\) there are \(A\), \(\bar a_0\) and
        \(\bar S_0\) with \(A\entails\chi(\bar a_0,\bar S_0)\) and
        \(\theta_0(A,\bar a_0,\bar S_0)\approx B\)
\end{enumerate}
\cite[p.~448, Definition~1.2.4]{Baldwin1985DSOQ}. The interpreting formulas
may contain a sequence \(\langle s_0,\dots,s_n\rangle\) of relation
variables, each ranging over \(\mathscr R_\phi(A)\). When \(T\) is the
theory of equality one writes \(Q_{\psi(r)}\le Q_{\phi(s)}\);
\(\equiv\) is interpretability in both directions (bi-interpretability).
For theories \(T_1,T_2\), \((T_1,Q_\psi)\le(T_2,Q_\phi)\) adds an
interpreting formula in the language of \((T_2,Q_\phi)\) for each relation
symbol of \(T_1\) \cite[p.~448]{Baldwin1985DSOQ}.
\end{hbdefinition}

\begin{hbconvention}{def:p03-w03-p03-046-b-reading}{Reading of Definition~1.2.4}
As printed, the coding formula \(\theta\) carries only the varying
parameters \(\bar a,\bar S\), so the family in clause~(i) cannot depend on
the guard parameters \(\bar a_0,\bar S_0\). As in
\cref{def:p03-w03-p03-046-reading}, \(\theta\) is here allowed to carry
\(\bar a_0,\bar S_0\) as well as \(\bar a,\bar S\). An interpretation is
\emph{full-domain} if \(\theta_0\) is \(x_0=x_0\). Baldwin notes that by
modifying \(\theta_0\) one may assume each structure
\((\theta_0(A),\theta(A))\) satisfies \(\psi\)
\cite[p.~448]{Baldwin1985DSOQ}; \cref{thm:p03-w03-p03-046-b-normal} proves
this remark.
\end{hbconvention}

The point of the remark is that the inclusion of
\(\mathscr R_\psi(\theta_0(A))\) in the coded family is the substantial half
of clause~(i): once every relation satisfying \(\psi\) on the domain is
coded, the coding can be repaired so that nothing else is.

\begin{hblemma}{thm:p03-w03-p03-046-b-normal}{proved=annex:p03-w03-p03-046-b-normal}{Normalizing the coded structures}
Let \(\psi(r)\), with \(r\) \(k\)-ary, hold of some relation on some
infinite set, and let \(\chi(\bar y_0,\bar s_0)\),
\(\theta_0(x_0,\bar y_0,\bar s_0)\) and
\(\theta(x_0,\dots,x_{k-1},\bar y_0,\bar s_0,\bar y,\bar s)\) be
first-order formulas satisfying clause~(ii) of Definition~1.2.4 and the
following half of clause~(i): if \(A\entails\chi(\bar a_0,\bar S_0)\), then
\(D=\theta_0(A,\bar a_0,\bar S_0)\) is infinite, \(\bar S_0\in\mathscr
R_\phi(A)\), and every \(R\in\mathscr R_\psi(D)\) equals
\(\theta(A,\bar a_0,\bar S_0,\bar a,\bar S)\cap D^k\) for some \(\bar a\in
A\) and \(\bar S\in\mathscr R_\phi(A)\). Then there are first-order formulas
\(\chi'\), \(\theta_0'\), \(\theta'\), with \(\theta_0'\) obtained from
\(\theta_0\) by adding a conjunct in new parameters, such that
\((\chi',\theta_0',\theta')\) satisfies Definition~1.2.4 and, for all
values of all parameters in any set \(A\) with
\(\theta_0'(A,\dots)\neq\varnothing\) and all \(\bar a\in A\), \(\bar
S\in\mathscr R_\phi(A)\), the structure
\((\theta_0'(A,\dots),\theta'(A,\dots,\bar a,\bar S))\) is a model of
\(\psi\).
\end{hblemma}

Baldwin attributes the modification to \(\theta_0\) alone; in the proof the
new conjunct of \(\theta_0'\) tests a default relation, and the guard and the
coding formula change with it.

In BS notation (\cref{def:p03-w03-p03-046-interpretation}) the guard
\(\chi\) of Definition~1.2.4 plays the role of \(\pi\), the domain formula
\(\theta_0\) that of \(\theta\), and the coding formula \(\theta\) that of
the \(\chi_i\).

Baldwin observes that \(Q_\psi\le Q_\phi\) implies
\((T,Q_\psi)\le(T,Q_\phi)\) for every theory \(T\), and, in the sense of
Chapter~II of the collection, \(L_{\omega\omega}(Q_\psi)\le
L_{\omega\omega}(Q_\phi)\) for every vocabulary, and likewise for
infinitary logics \cite[p.~448]{Baldwin1985DSOQ}. The form proved here is
the following; it covers every interpretation constructed in this
section, all of which are full-domain, and every interpretation in
\(Q_{\mathrm I}\) (\cref{thm:p03-w03-p03-046-b-qi-normal}). Since
\(\mathscr R_\psi(A)\) is defined only for infinite \(A\) (Notation~1.2.1),
the comparison is made on infinite structures.

\begin{hblemma}{thm:p03-w03-p03-046-b-transfer}{proved=annex:p03-w03-p03-046-b-transfer}{Transfer along a full-domain interpretation}
Let \(\mathcal Q\) be a set of definable second-order quantifiers containing
\(Q_\phi\), and suppose \(Q_\psi\le Q_\phi\) by a full-domain
interpretation \((\chi,x_0=x_0,\theta)\).
\begin{enumerate}
  \item For every vocabulary \(L\), all infinite cardinals
        \(\kappa,\lambda\) and every formula \(\sigma\) of
        \(L_{\kappa\lambda}(\mathcal Q\cup\{Q_\psi\})\) (conjunctions and
        disjunctions of fewer than \(\kappa\) formulas, blocks of fewer than
        \(\lambda\) first-order quantifiers) there is a formula \(\sigma'\)
        of \(L_{\kappa\lambda}(\mathcal Q)\) with the same free variables
        such that \(M\entails\sigma\leftrightarrow\sigma'\) under every
        assignment in every infinite \(L\)-structure \(M\). The case
        \(\kappa=\lambda=\omega\) is the finitary logic.
  \item For every first-order theory \(T\) all of whose models are
        infinite, \((T,Q_\psi)\le(T,Q_\phi)\) in the extended sense of
        Definition~1.2.4, by the guard \(\chi\), the domain formula
        \(x_0=x_0\), the coding formula \(\theta\), and each relation symbol
        of \(L(T)\) interpreted by itself.
\end{enumerate}
\end{hblemma}

The full-domain hypothesis matters: if \(\theta_0\) defines a proper subset,
the translation of \(Q_\psi\) acts on that subset and not on \(M\), and a
relation of an \(L\)-structure on \(M\) cannot in general be carried to the
subset. Baldwin's observation asserts the existence of an interpretation
\((T,Q_\psi)\le(T,Q_\phi)\), not that the given one works, so the question is
whether every interpretation can be replaced by a full-domain one.
\Cref{thm:p03-w03-p03-046-b-qi-normal,thm:p03-w03-p03-046-b-full-domain}
below answer this for interpretations in \(Q_{\mathrm I}\) and for every
quantifier \(Q_\phi\) that interprets \(Q_{1-1}\) with full domain. For
interpretations that are not full-domain a weaker transfer always holds,
which is the form in which Hanf-number arguments use them.

\begin{hblemma}{thm:p03-w03-p03-046-b-relativize}{proved=annex:p03-w03-p03-046-b-relativize}{Relativized transfer}
Suppose \(Q_\psi\le Q_\phi\) by \((\chi,\theta_0,\theta)\). For every
\(Q_\psi\)-sentence \(\sigma\) of pure identity theory there is a
\(Q_\phi\)-sentence \(\sigma^\sharp\) of pure identity theory such that,
for every infinite set \(A\), \(A\entails\sigma^\sharp\) iff there are
\(\bar a_0\in A\) and \(\bar S_0\) with \(A\entails\chi(\bar a_0,\bar S_0)\)
and \(\theta_0(A,\bar a_0,\bar S_0)\entails\sigma\). If \(\sigma\) has an
infinite model, then \(\sigma^\sharp\) has one.
\end{hblemma}

Interpretations in first-order quantification can always be replaced by
full-domain ones, because a relation coded from finitely many points of a
pure set is fixed by the equality pattern of those points.

\begin{hblemma}{thm:p03-w03-p03-046-b-qi-normal}{proved=annex:p03-w03-p03-046-b-qi-normal}{Interpretations in \(Q_{\mathrm I}\) are full-domain}
Suppose \(Q_{\psi(r)}\le Q_{\mathrm I}\), with \(r\) \(k\)-ary. Then:
\begin{enumerate}
  \item there are finitely many first-order formulas
        \(\delta_1(\bar x,\bar z),\dots,\delta_N(\bar x,\bar z)\) of pure
        identity theory, \(\bar x=\langle x_0,\dots,x_{k-1}\rangle\), such
        that for every infinite set \(M\) every \(R\in\mathscr R_\psi(M)\)
        equals \(\delta_i(M,\bar c)\) for some \(i\) and some \(\bar c\in M\);
  \item \(Q_\psi\le Q_{\mathrm I}\) by a full-domain interpretation.
\end{enumerate}
\end{hblemma}

Baldwin's remark that interpretations preserve Hanf number
\cite[p.~448]{Baldwin1985DSOQ} uses the Hanf number of \((T,Q)\) in the
following sense: the least cardinal \(\mu\) such that every sentence of
\((T,Q)\) with a model of cardinality at least \(\mu\) has arbitrarily large
models. Here \(h(Q)\) is the Hanf number of \((\mathrm{Th}(=),Q)\) and
\(\ell(Q)\) the least \(\mu\) such that every \(Q\)-sentence of pure identity
theory with an infinite model has one of cardinality at most \(\mu\).

\begin{hbcorollary}{thm:p03-w03-p03-046-b-hanf}{proved=annex:p03-w03-p03-046-b-hanf}{Interpretations preserve Hanf numbers}
Suppose \(Q_\psi\le Q_\phi\) by a full-domain interpretation, or
\(Q_\psi\le Q_{\mathrm I}\) by any interpretation. Then every
\(Q_\psi\)-sentence of pure identity theory has the same infinite models as
a \(Q_\phi\)-sentence, and consequently \(h(Q_\psi)\le h(Q_\phi)\) and
\(\ell(Q_\psi)\le\ell(Q_\phi)\) whenever the right-hand numbers exist.
\end{hbcorollary}

For an interpretation that is not full-domain, the domain
\(\theta_0(A,\bar a_0,\bar S_0)\) may be much smaller than \(A\): the
relativized sentence \(\sigma^\sharp\) of
\cref{thm:p03-w03-p03-046-b-relativize} then has arbitrarily large models
whenever \(\sigma\) has a large one, but nothing forces the domains inside
those large models to be large. Baldwin returns to this point in his
Section~4, where an interpretation is required to compute the size of the
interpreted model from that of the interpreting one
\cite[pp.~470--471]{Baldwin1985DSOQ}
(\cref{def:p03-w03-p03-046-b-semantic}). When the interpreting quantifier
can code permutations, the domain can be moved onto the whole set.

\begin{hblemma}{thm:p03-w03-p03-046-b-full-domain}{proved=annex:p03-w03-p03-046-b-full-domain}{Full domains when \(Q_{1-1}\) is interpreted with full domain}
Suppose \(Q_{1-1}\le Q_\phi\) by a full-domain interpretation. Then:
\begin{enumerate}
  \item for every infinite set \(A\) and every interpretation
        \((\chi,\theta_0,\theta)\) of a quantifier \(Q_\psi\) in \(Q_\phi\)
        there are guard parameters \(\bar a_0,\bar S_0\) in \(A\) with
        \(|\theta_0(A,\bar a_0,\bar S_0)|=|A|\);
  \item if \(Q_\psi\le Q_\phi\) by any interpretation, then
        \(Q_\psi\le Q_\phi\) by a full-domain interpretation.
\end{enumerate}
Part~(1) holds for every interpretation in every \(Q_\phi\), without the
hypothesis on \(Q_{1-1}\).
\end{hblemma}

\begin{hbcorollary}{thm:p03-w03-p03-046-b-transfer-general}{proved=annex:p03-w03-p03-046-b-transfer-general}{Transfer and Hanf numbers for quantifiers coding pairings}
Let \(Q_\phi\) be \(Q_{1-1}\), \(Q_{\mathrm{inv}}\), \(Q_{\mathrm{perm}}\),
\(Q_{\mathrm{II}}\), or any quantifier into which \(Q_{1-1}\) is
interpretable with full domain, for instance one satisfying the hypothesis
of \cref{thm:p03-w03-p03-046-b-perm-crit}. If \(Q_\psi\le Q_\phi\) by any
interpretation, then the conclusions of
\cref{thm:p03-w03-p03-046-b-transfer}(1) and~(2) hold: for every theory
\(T\) all of whose models are infinite, \((T,Q_\psi)\le(T,Q_\phi)\), and
\(L_{\kappa\lambda}(Q_\psi)\)-definable classes of infinite structures are
\(L_{\kappa\lambda}(Q_\phi)\)-definable in every vocabulary; moreover
\(h(Q_\psi)\le h(Q_\phi)\) and \(\ell(Q_\psi)\le\ell(Q_\phi)\) whenever the
right-hand numbers exist.
\end{hbcorollary}

Together with \cref{thm:p03-w03-p03-046-b-qi-normal}, Baldwin's p.~448
claims are thus proved for every interpretation whose target is
\(Q_{\mathrm I}\) or a quantifier that interprets \(Q_{1-1}\) with full
domain. For a target \(Q_\phi\) that interprets neither, the argument
breaks at one point: the subsets that \(Q_\phi\) supplies do not give a
bijection between \(A\) and the domain \(\theta_0(A,\bar a_0,\bar S_0)\),
and for such targets the claims for interpretations that are not
full-domain are not proved here. By the four-quantifier classification of
Baldwin's Section~2 (\cref{def:p03-w03-p03-046-only-four}) such a target is
bi-interpretable with \(Q_{\mathrm I}\) or \(Q_{\mathrm{mon}}\); that
classification is not proved in this chapter.

\subsection{Simple definability}
\label{sec:p03-w03-p03-046-b-simple}

\begin{hbdefinition}{def:p03-w03-p03-046-b-simply}{Simple definability}
If the relation \(S\) on \(A\) is defined by \(\phi(\bar x,\bar b,\bar R)\)
with \(\bar R\in\mathscr R_\psi(A)\), where \(\phi\) is first-order, then
\(S\) is \emph{simply definable by \(Q_\psi\)}
\cite[p.~449, Definition~1.2.6]{Baldwin1985DSOQ}.
\end{hbdefinition}

\begin{hblemma}{thm:p03-w03-p03-046-b-maximal}{proved=annex:p03-w03-p03-046-b-maximal}{\(Q_{\mathrm{II}}\) is maximal (Proposition~1.2.7)}
For every definable second-order quantifier \(Q_\psi\) whose condition
\(\psi\) holds of some relations on some infinite set,
\(Q_\psi\le Q_{\mathrm{II}}\), by a full-domain interpretation
\cite[p.~449, Proposition~1.2.7]{Baldwin1985DSOQ}.
\end{hblemma}

Baldwin states Proposition~1.2.7 for every \(\psi\). If \(\psi\) holds of no
relations on any infinite set, then \(\mathscr R_\psi(B)=\varnothing\) for
every infinite \(B\), and clause~(i) of Definition~1.2.4 would require
the family \(\{\theta(A,\bar a,\bar S)\}\), indexed by a nonempty set of
parameters, to be empty; so such a \(Q_\psi\), whose quantified formulas
are all false, is interpretable in no definable quantifier in the sense of
that definition. By compactness (\cref{thm:p03-compactness-ultraproduct}) and the
downward L\"owenheim--Skolem theorem
(\cref{thm:p03-w03-p03-039-downward-ls}, for the fragment of first-order
formulas), \(\psi\) holds of relations on one infinite set iff it does on
every infinite set.

\begin{hblemma}{thm:p03-w03-p03-046-b-mon}{proved=annex:p03-w03-p03-046-b-mon}{An infinite, coinfinite set gives monadic logic (Lemma~1.2.8)}
If \(Q_\psi\) simply defines an infinite, coinfinite set, then
\(Q_{\mathrm{mon}}\le Q_\psi\) \cite[p.~449, Lemma~1.2.8]{Baldwin1985DSOQ}.
More precisely, there is a first-order formula \(\mu(x,\bar z,\bar r)\)
such that on every infinite set \(A\) the sets \(\mu(A,\bar b,\bar R)\),
\(\bar b\in A\), \(\bar R\in\mathscr R_\psi(A)\), are exactly the subsets
of \(A\), and the interpretation is full-domain. Conversely, if
\(Q_{\mathrm{mon}}\le Q_\psi\), then \(Q_\psi\) simply defines an infinite,
coinfinite set.
\end{hblemma}

Baldwin's proof calls \(X\subseteq A\) \emph{regular} if
\(|X|=|A-X|=|A|\), transports one definable regular set to every other by
a permutation of \(A\), which preserves \(\psi\) because \(\psi\) has no
nonlogical symbols, and writes every subset as a Boolean combination of
regular sets \cite[p.~449]{Baldwin1985DSOQ}. The annex proof follows this
route and adds the compactness and L\"owenheim--Skolem step that produces
a definable regular set in every infinite cardinality from an infinite,
coinfinite one.

Read contrapositively, \cref{thm:p03-w03-p03-046-b-mon} says that when
\(Q_{\mathrm{mon}}\not\le Q_\psi\) every infinite model \((A,\bar R)\) of
\(T_\psi\) is \emph{strongly minimal}: every subset definable with
parameters is finite or cofinite, in the sense of Baldwin and Lachlan
\cite[p.~449]{Baldwin1985DSOQ}. Baldwin adds that compactness then bounds
finite definable families uniformly \cite[pp.~449--450]{Baldwin1985DSOQ}.

\begin{hblemma}{thm:p03-w03-p03-046-b-uniform}{proved=annex:p03-w03-p03-046-b-uniform}{Strong minimality and uniform bounds}
Suppose \(Q_{\mathrm{mon}}\not\le Q_\psi\). Then for every infinite \(A\)
and \(\bar R\in\mathscr R_\psi(A)\) the structure \((A,\bar R)\) is
strongly minimal, and for every first-order \(\phi(x,\bar y,\bar r)\), if
\(\phi(A,\bar b,\bar R)\) is finite for every \(\bar b\in A\), then there
is \(k<\omega\) with \(|\phi(A,\bar b,\bar R)|\le k\) for all \(\bar b\).
\end{hblemma}

\begin{hblemma}{thm:p03-w03-p03-046-b-perm-crit}{proved=annex:p03-w03-p03-046-b-perm-crit}{An equivalence relation gives permutational logic (Lemma~1.2.9)}
If \(Q_\psi\) simply defines an equivalence relation with infinitely many
infinite classes and \(Q_{\mathrm{mon}}\le Q_\psi\), then
\(Q_{1-1}\le Q_\psi\), by a full-domain interpretation
\cite[p.~450, Lemma~1.2.9]{Baldwin1985DSOQ}.
\end{hblemma}

Baldwin's proof uses only infinitely many classes with at least two
elements on some infinite subset, shrinks to classes of exactly two
elements with an infinite complement, and then uses compactness, the
L\"owenheim--Skolem theorems and the absence of nonlogical symbols in
\(\psi\) to spread the pairing over every infinite set
\cite[p.~450]{Baldwin1985DSOQ}. The annex proof carries out these steps and
adds the default clause that makes the family of coded relations exactly
the pairings.

\subsection{Permutations, involutions and pairings}
\label{sec:p03-w03-p03-046-b-permutations}

Baldwin explains the passage from arbitrary permutations to permutations
of order two by replacing a permutation \(f\) by three involutions
\(g,h,j\): on each orbit of \(f\), \(g\) fixes every other element, and
\(h\) and \(j\) are products of two-cycles agreeing with \(f\) on the
points fixed by \(g\) and on the points moved by \(g\), respectively
\cite[p.~445]{Baldwin1985DSOQ}. This description covers the infinite
orbits and the orbits of length at most two, but not a finite orbit of odd
length at least three: if \(x\) and \(f(x)\) were both fixed by \(g\), then
\(h\) would exchange \(x\) and \(f(x)\) and also \(f(x)\) and \(f(f(x))\),
forcing \(f(f(x))=x\), and the same holds with \(j\) for two consecutive
points moved by \(g\); so on such an orbit the fixed and the moved points
of \(g\) would alternate around a cycle of odd length. The proof below uses
instead the factorization of every permutation into two involutions.

\begin{hbtheorem}{thm:p03-w03-p03-046-b-permutations}{proved=annex:p03-w03-p03-046-b-permutations}{Permutations, involutions and pairings}
Let \(Q_{\mathrm{perm}}\) quantify over graphs of permutations and
\(Q_{\mathrm{inv}}\) over graphs of involutions (permutations \(\alpha\)
with \(\alpha\circ\alpha=\mathrm{id}\), fixed points allowed); both are
definable second-order quantifiers with a binary variable. Then
\(Q_{\mathrm{perm}}\le Q_{\mathrm{inv}}\le Q_{1-1}\) and
\(Q_{1-1}\le Q_{\mathrm{inv}}\le Q_{\mathrm{perm}}\), all by full-domain
interpretations. In particular every permutation of a set is the
composite of two involutions.
\end{hbtheorem}

This is the content of Baldwin's remarks that quantification over
arbitrary permutations is bi-interpretable with quantification over
permutations of order two, and the latter with the permutational
quantifier \cite[pp.~445, 449]{Baldwin1985DSOQ}. It also bridges the two
notations: by \cref{thm:p03-w03-p03-046-b-transfer}, for every theory \(T\)
the logic \(\text{1-1}\) of \cref{def:p03-w03-p03-046-four} and
Baldwin's \(Q_{1-1}\) define the same classes of infinite models.

\subsection{The chain and its separation}
\label{sec:p03-w03-p03-046-b-chain}

\begin{hblemma}{thm:p03-w03-p03-046-b-mso-sets}{proved=annex:p03-w03-p03-046-b-mso-sets}{Monadic logic of unary structures}
Let \(L\) consist of finitely many unary predicates (possibly none).
\begin{enumerate}
  \item Every sentence of monadic second-order logic in \(L\) that has an
        infinite model has a countable model and models of every infinite
        cardinality; so for this vocabulary monadic logic has L\"owenheim
        and Hanf number \(\aleph_0\).
  \item Every sentence of \((\mathrm{Th}(=),Q_{\mathrm{mon}})\) is either
        true on all infinite sets or false on all infinite sets; so the
        Hanf number of \((\mathrm{Th}(=),Q_{\mathrm{mon}})\) is \(\aleph_0\).
\end{enumerate}
\end{hblemma}

Both statements are Baldwin's \cite[pp.~446, 449]{Baldwin1985DSOQ}; the
chapter does not prove them, and the annex proves them by a
cardinality-threshold argument on the cells of the Boolean algebra
generated by the predicates and set variables.

Baldwin also says that for a vocabulary of unary predicates the
Feferman--Vaught theorem holds for monadic logic
\cite[p.~446]{Baldwin1985DSOQ}, without stating a form. The form proved
here covers the two operations on which Feferman and Vaught's theorem is
modelled, disjoint unions and direct products, with the reduction made
explicit. For \(L=\{P_1,\dots,P_m\}\) and \(L\)-structures
\(\mathfrak A,\mathfrak B\), the \emph{disjoint union}
\(\mathfrak A\sqcup\mathfrak B\) has as universe the disjoint union of \(A\)
and \(B\) and interprets each \(P_j\) as the union of \(P_j^{\mathfrak A}\)
and \(P_j^{\mathfrak B}\); the \emph{direct product}
\(\mathfrak A\times\mathfrak B\) has universe \(A\times B\) and
\(P_j^{\mathfrak A\times\mathfrak B}=P_j^{\mathfrak A}\times
P_j^{\mathfrak B}\). The quantifier rank of a monadic formula counts
individual and set quantifiers, and \(t_0=0\), \(t_{n+1}=2t_n+1\).

\begin{hblemma}{thm:p03-w03-p03-046-b-fv}{proved=annex:p03-w03-p03-046-b-fv}{A Feferman--Vaught theorem for unary monadic logic}
Let \(L=\{P_1,\dots,P_m\}\) and \(n<\omega\). There are finitely many
first-order \(L\)-sentences \(\delta_1,\dots,\delta_N\), each structure
satisfying exactly one of them, such that for every sentence \(\sigma\) of
monadic second-order logic in \(L\) of quantifier rank at most \(n\) there
are sets \(\Gamma^\sqcup_\sigma,\Gamma^\times_\sigma\subseteq\{1,\dots,N\}^2\)
with, for all \(L\)-structures \(\mathfrak A,\mathfrak B\),
\[
\mathfrak A\sqcup\mathfrak B\entails\sigma\iff
\bigvee_{(i,j)\in\Gamma^\sqcup_\sigma}\bigl(\mathfrak A\entails\delta_i
\text{ and }\mathfrak B\entails\delta_j\bigr),
\qquad
\mathfrak A\times\mathfrak B\entails\sigma\iff
\bigvee_{(i,j)\in\Gamma^\times_\sigma}\bigl(\mathfrak A\entails\delta_i
\text{ and }\mathfrak B\entails\delta_j\bigr).
\]
The sentences \(\delta_i\) have quantifier rank at most \(t_n\). In
particular, if \(\mathfrak A\) and \(\mathfrak A'\) satisfy the same
monadic sentences, and so do \(\mathfrak B\) and \(\mathfrak B'\), then
\(\mathfrak A\sqcup\mathfrak B\) and \(\mathfrak A'\sqcup\mathfrak B'\)
satisfy the same monadic sentences, and so do \(\mathfrak A\times\mathfrak
B\) and \(\mathfrak A'\times\mathfrak B'\).
\end{hblemma}

Baldwin's Section~3 develops a Feferman--Vaught theorem for infinitary
monadic logic and abstractly defined generalized products, of which the
disjoint union of two structures is a case \cite[pp.~466, 468]{Baldwin1985DSOQ};
that theorem is not presented in this chapter.

\begin{hblemma}{thm:p03-w03-p03-046-b-uncountable}{proved=annex:p03-w03-p03-046-b-uncountable}{A permutational sentence with only uncountable models}
There is a sentence of \((\mathrm{Th}(=),Q_{1-1})\) whose infinite models
are exactly the uncountable sets.
\end{hblemma}

Baldwin obtains such sentences from Isbell's theorem that some sentence of
group theory holds in the symmetric group \(S_\kappa\) exactly when
\(\kappa=\omega\) \cite[pp.~445, 449]{Baldwin1985DSOQ}; Isbell's proof is not
in the chapter. The sentence of the annex proof says instead that no
permutation of the universe has a single orbit, using monadic
quantification inside \(Q_{1-1}\).

\begin{hbtheorem}{thm:p03-w03-p03-046-b-distinct}{proved=annex:p03-w03-p03-046-b-distinct}{The chain \(Q_{\mathrm I}<Q_{\mathrm{mon}}<Q_{1-1}\le Q_{\mathrm{II}}\) (Theorem~1.2.5, first two parts)}
\(Q_{\mathrm I}\le Q_{\mathrm{mon}}\le Q_{1-1}\le Q_{\mathrm{II}}\), and
\(Q_{\mathrm{mon}}\not\le Q_{\mathrm I}\) and
\(Q_{1-1}\not\le Q_{\mathrm{mon}}\)
\cite[p.~449, Theorem~1.2.5]{Baldwin1985DSOQ}.
\end{hbtheorem}

Baldwin separates \(Q_{\mathrm{mon}}\) from \(Q_{\mathrm I}\) by noting that
the class of well-orders is definable in \((\mathrm{Th}(<),Q_{\mathrm{mon}})\)
but not in \((\mathrm{Th}(<),Q_{\mathrm I})\), and appeals to the transfer
observation \cite[p.~449]{Baldwin1985DSOQ}. The well-order claim is the
following proposition, and the transfer it needs is available because every
interpretation in \(Q_{\mathrm I}\) can be made full-domain
(\cref{thm:p03-w03-p03-046-b-qi-normal,thm:p03-w03-p03-046-b-transfer}).

\begin{hblemma}{thm:p03-w03-p03-046-b-wellorder}{proved=annex:p03-w03-p03-046-b-wellorder}{Well-orders in monadic and first-order logic of order}
\begin{enumerate}
  \item The sentence
        \(\mathrm{WO}:=\neg Q_{\mathrm{mon}}r\,\bigl(\exists x\,r(x)\wedge
        \forall x\,(r(x)\to\exists y\,(r(y)\wedge y<x))\bigr)\) of
        \((\mathrm{Th}(<),Q_{\mathrm{mon}})\) holds in a linear order
        \((M,<)\) iff \((M,<)\) is a well-order.
  \item No sentence of \((\mathrm{Th}(<),Q_{\mathrm I})\) holds in exactly
        the infinite well-orders among the infinite linear orders.
\end{enumerate}
\end{hblemma}

The annex proves \(Q_{\mathrm{mon}}\not\le Q_{\mathrm I}\) along this route
and also by counting definable sets in a countable model, which uses no
transfer. The separation of \(Q_{1-1}\) from
\(Q_{\mathrm{mon}}\) follows Baldwin's route: the monadic theory of
infinite sets is constant (\cref{thm:p03-w03-p03-046-b-mso-sets}), while
\(Q_{1-1}\) distinguishes countable from uncountable sets
(\cref{thm:p03-w03-p03-046-b-uncountable}).

\begin{proofboundary}{def:p03-w03-p03-046-b-third}{Third part of Theorem~1.2.5}
Baldwin completes Theorem~1.2.5 with \(Q_{\mathrm{II}}\not\le Q_{1-1}\),
deriving it from Shelah's computation that the Hanf number of
\((\mathrm{Th}(=),Q_{1-1})\) is \(\aleph_{\Omega^\omega}\), where
\(\Omega=(2^\omega)^+\) \cite[pp.~445, 449]{Baldwin1985DSOQ}. Shelah's
computation is not reproduced in the chapter and is not proved here, so
the strictness of the last link of the chain is not established in this
section.
\end{proofboundary}

The main result, proved in Section~2 of the chapter, is that up to
bi-interpretability the four quantifiers are the only definable
second-order quantifiers and are linearly ordered by interpretation, and
that quantifiers with finite strings of relation variables give no new
cases \cite[p.~450]{Baldwin1985DSOQ}; the BS form of the classification
is recorded in \cref{def:p03-w03-p03-046-only-four}. Two of the four may
coincide over a particular theory: in the presence of a pairing function
\(Q_{\mathrm{mon}}\) is equivalent to \(Q_{\mathrm{II}}\)
\cite[p.~448]{Baldwin1985DSOQ} (\cref{thm:p03-w03-p03-046-pairing}).

\subsection{Conditions for interpretability}
\label{sec:p03-w03-p03-046-b-conditions}

\begin{hbdefinition}{def:p03-w03-p03-046-b-codable}{Codable theory}
The theory \(T\) is \emph{codable} if, for some \(n\) and some model \(M\)
of \(T\), there are infinite sets \(\langle B_i:i<n\rangle\) and \(C\)
contained in \(M\) and a first-order formula, possibly with parameters,
\(\phi(x,y_0,\dots,y_{n-1})\), which defines a one-to-one map from
\(B_0\times\cdots\times B_{n-1}\) onto \(C\)
\cite[p.~450, Definition~1.3.1]{Baldwin1985DSOQ}. Here \(n\ge2\) and
``defines'' means: for each \(\bar b\in B_0\times\cdots\times B_{n-1}\)
there is exactly one \(x\in M\) with \(M\entails\phi(x,\bar b)\), this
\(x\) lies in \(C\), distinct tuples give distinct \(x\), and every member
of \(C\) is obtained.
\end{hbdefinition}

The restriction \(n\ge2\) is forced: with \(n=1\) every theory with an
infinite model would be codable, by \(\phi(x,y_0):=x=y_0\) and
\(B_0=C\).

\begin{hbtheorem}{thm:p03-w03-p03-046-b-codable}{proved=annex:p03-w03-p03-046-b-codable}{Codable theories (Theorem~1.3.2)}
If \(T\) is codable, then \(Q_{\mathrm{II}}\le(T,Q_{\mathrm{mon}})\)
\cite[p.~450, Theorem~1.3.2]{Baldwin1985DSOQ}.
\end{hbtheorem}

Baldwin explains the theorem by the pairing function from two sets of
power \(\kappa\) onto a third that codability yields in every cardinality,
with which a binary relation is coded by a subset of the third set, and
refers to Section~II.2.4 of Baldwin and Shelah for the details
\cite[p.~450]{Baldwin1985DSOQ}; the annex proof reduces the theorem to
\cref{thm:p03-w03-p03-046-coding}.

\begin{hbtheorem}{thm:p03-w03-p03-046-b-equiv}{proved=annex:p03-w03-p03-046-b-equiv}{Equivalence relations with infinitely many infinite classes (Lemma~1.3.3)}
If there is a first-order formula \(\phi(x,y)\) which defines, on some
model \(M\) of a first-order theory \(T\) and on some infinite subset
\(A\) of \(M\), an equivalence relation with infinitely many infinite
classes, then \(Q_{\mathrm{II}}\le(T,Q_{1-1})\)
\cite[p.~451, Lemma~1.3.3]{Baldwin1985DSOQ}.
\end{hbtheorem}

Baldwin refers this lemma to Section~II.2.6 of Baldwin and Shelah
\cite[p.~451]{Baldwin1985DSOQ}; the annex proof combines
\cref{thm:p03-w03-p03-046-equiv} with
\cref{thm:p03-w03-p03-046-b-permutations}.

\begin{hbexample}{def:p03-w03-p03-046-b-example}{Simple definability by pairings on a countable set}
Let \(A=\omega\). Let \(E_1\) pair \(2i\) with \(2i+1\) for every \(i\).
Let \(E_2\) pair \(4i\) with \(4i+1\) for every \(i\), and, on each block
\(\{8j+2,8j+3,8j+6,8j+7\}\), pair \(8j+2\) with \(8j+6\) and \(8j+3\) with
\(8j+7\). Both are pairings of \(\omega\), and the formula
\(\phi(x,E_1,E_2):=\forall y\,(E_1(x,y)\leftrightarrow E_2(x,y))\)
defines the set \(X_0=\{4i,4i+1:i<\omega\}\) of points with the same
partner under both, which is infinite and coinfinite. So \(Q_{1-1}\)
simply defines an infinite, coinfinite set and
\cref{thm:p03-w03-p03-046-b-mon} applies.
\begin{enumerate}
  \item \(X_0\) is regular in \(\omega\). The set \(X=3\omega\) of
        multiples of \(3\) is regular too, so some permutation \(g\) of
        \(\omega\) maps \(X_0\) onto \(X\) and \(\omega-X_0\) onto
        \(\omega-X\); the pairings \(g[E_1]\), \(g[E_2]\) satisfy
        \(\phi(\omega,g[E_1],g[E_2])=X\).
  \item The finite set \(\{5\}\) is \(Y\cap(\omega-Z)\) with \(Y=\{5\}\cup
        2\omega\) and \(Z=2\omega\), both regular; each of \(Y\), \(Z\) is
        \(\phi(\omega,E,E')\) for suitable pairings, by the same transport.
  \item The pairing \(E_1\) itself is obtained from any other pairing
        \(E\) of \(\omega\) by a permutation sending the blocks of \(E\) to
        those of \(E_1\), which is the step that makes the family of
        coded pairings in \cref{thm:p03-w03-p03-046-b-perm-crit} exhaustive.
\end{enumerate}
\end{hbexample}

\begin{hbwarning}{def:p03-w03-p03-046-b-limits}{Limits of the results of this section}
\begin{enumerate}
  \item The transfer from \(Q_\psi\le Q_\phi\) to
        \((T,Q_\psi)\le(T,Q_\phi)\), and the preservation of Hanf numbers,
        are proved for full-domain interpretations, for interpretations in
        \(Q_{\mathrm I}\), and for interpretations in any quantifier that
        interprets \(Q_{1-1}\) with full domain
        (\cref{thm:p03-w03-p03-046-b-transfer-general}); for other targets
        only the relativized form \cref{thm:p03-w03-p03-046-b-relativize} is
        available.
  \item The strictness of \(Q_{1-1}\le Q_{\mathrm{II}}\) rests on Shelah's
        Hanf-number computation, which is reported, not proved
        (\cref{def:p03-w03-p03-046-b-third}).
  \item A quantifier condition with no infinite models is interpretable in
        no definable quantifier under Definition~1.2.4, although its
        quantified formulas are all false.
  \item The three-involution description of a permutation fails on finite
        orbits of odd length at least three; the two-involution
        factorization does not.
\end{enumerate}
\end{hbwarning}

\section{Comparison of theories and further quantifiers}
\label{sec:p03-w03-p03-046-b-comparison}

This section follows Sections~4 and~6 of Baldwin's chapter
\cite[pp.~470--472, 476--477]{Baldwin1985DSOQ}. Section~4 separates the
notion of interpretation that preserves Hanf numbers from a purely
syntactic one; Section~6 lists extensions of the notion of definable
quantifier studied by Shelah.

\emph{Purpose.} Baldwin's point is that a proof of non-interpretability by
a computation of Hanf numbers needs an interpretation that controls
cardinalities, and that this is a real restriction: Gurevich and Shelah
showed it consistent, indeed a consequence of the generalized continuum
hypothesis, that \(Q_{\mathrm{II}}\) has a syntactic interpretation in the
monadic theory of order \cite[p.~470]{Baldwin1985DSOQ}, while Baldwin's
Theorem~4.6 excludes a strong semantic one.

\begin{hbdefinition}{def:p03-w03-p03-046-b-syntactic}{Syntactic interpretability (Definition~4.1)}
A theory \(T_1\) is \emph{syntactically interpretable} in the logic
\(T_2\) if there is a map \(*\) assigning to each \(T_1\)-sentence \(\phi\)
a \(T_2\)-sentence \(\phi^*\) such that \(T_1\vdash\phi\) iff
\(T_2\vdash\phi^*\) \cite[p.~470, Definition~4.1]{Baldwin1985DSOQ}. Here a
theory in a logic with second-order quantifiers is a set of sentences
\((T,Q_\psi)\) as in \cref{def:p03-w03-p03-046-b-quantifier}, and
\(T_i\vdash\phi\) is read as \(\phi\in T_i\), that is, \(\phi\) holds in
every model of the first-order theory underlying \(T_i\). This is
Definition~1.4.1 of Baldwin and Shelah
(\cref{def:p03-w03-p03-046-syntactic}).
\end{hbdefinition}

\begin{hblemma}{thm:p03-w03-p03-046-b-turing}{proved=annex:p03-w03-p03-046-b-turing}{Recursive syntactic interpretations and Turing degrees}
Suppose the sentences of both logics are coded by natural numbers so that
the codes of sentences form recursive sets, and the first-order theory
underlying \(T_1\) has a model. If \(T_1\) is syntactically interpretable
in \(T_2\) by a map \(*\) that is recursive on codes, then the set of
\(T_1\)-sentences true in \(T_1\) is many-one reducible, and hence Turing
reducible, to the corresponding set for \(T_2\); so the Turing degree of
\(T_1\) is at most that of \(T_2\) \cite[p.~470]{Baldwin1985DSOQ}.
\end{hblemma}

Baldwin stresses that such a map need not preserve model-theoretic
properties \cite[p.~470]{Baldwin1985DSOQ}: the Gurevich--Shelah syntactic
interpretation of \(Q_{\mathrm{II}}\) in \((\mathrm{Th}(<),Q_{\mathrm{mon}})\)
coexists with a gap between the Hanf numbers of the two logics. The
Gurevich--Shelah consistency result is reported from the chapter, which
does not prove it, and it is not proved here.

\begin{hbdefinition}{def:p03-w03-p03-046-b-semantic}{Semantic and strong semantic interpretation (Definition~4.2)}
The logic \(T_1\) is \emph{semantically interpretable} in the logic \(T_2\)
if there are maps, both written \(*\), sending each \(T_1\)-sentence
\(\phi\) to a \(T_2\)-sentence \(\phi^*\) and each model \(M\) of \(T_1\) to
a model \(M^*\) of \(T_2\), such that (i) \(M\subseteq M^*\) and (ii)
\(M\entails\phi\) iff \(M^*\entails\phi^*\). If in addition (iii)
\(|M^*|\) can be computed from \(|M|\), \(T_1\) is \emph{strongly
semantically interpretable} in \(T_2\)
\cite[p.~471, Definition~4.2]{Baldwin1985DSOQ}.
\end{hbdefinition}

The printed definition says that the second map takes the models of
\(T_2\) onto the models of \(T_2\) \cite[p.~471]{Baldwin1985DSOQ}; since
clause~(ii) evaluates the \(T_1\)-sentence \(\phi\) in \(M\), the domain of
the map is read here as the class of models of \(T_1\). Clause~(iii) is
read as: there is a function \(g\) on cardinals with \(|M^*|=g(|M|)\) for
every model \(M\) of \(T_1\).

\begin{hbdefinition}{def:p03-w03-p03-046-b-bounded}{Bounded Hanf numbers (Definition~4.3, Notation~4.4)}
Write \(H(T)\) for the Hanf number of a theory \(T\). The Hanf number of
\(T_1\) is \emph{bounded in terms of} that of \(T_2\), written
\(B(T_1,T_2)\), if there is a second-order definable function \(f(x)\) with
\(H(T_1)\le f(H(T_2))\) \cite[p.~471, Definition~4.3]{Baldwin1985DSOQ}.
\(HL\), \(HW\) and \(H_{\mathrm{II}}\) denote the Hanf numbers of
\((\mathrm{Th}(<),Q_{\mathrm{mon}})\), of the monadic theory of
well-orderings, and of \(Q_{\mathrm{II}}\), respectively
\cite[p.~471, Notation~4.4]{Baldwin1985DSOQ}.
\end{hbdefinition}

The chapter does not say what makes a function on cardinals second-order
definable. Baldwin uses the relation \(B\) for three claims: that it is
transitive; that \(B(Q_{\mathrm{II}},T)\) excludes a strong semantic
interpretation of \(Q_{\mathrm{II}}\) in \(T\); and, through his Lemma~4.5
(\(HL\le\sum\{2^\lambda:\lambda<HW\}\), proved from the decomposition of a
linear order as a free union of intervals and the preservation theorem of
his Section~3), that \(Q_{\mathrm{II}}\) has no strong semantic
interpretation in \((\mathrm{Th}(<),Q_{\mathrm{mon}})\), his Theorem~4.6
\cite[pp.~471--472]{Baldwin1985DSOQ}. These results are not proved in
this chapter.

Baldwin's Section~6 records ways of widening the notion of definable
second-order quantifier of \cref{def:p03-w03-p03-046-b-quantifier}
\cite[pp.~476--477]{Baldwin1985DSOQ}:
\begin{enumerate}
  \item replacing the first-order quantifier condition \(\psi\) by a
        formula of another logic, giving quantifiers definable in an
        infinitary logic, in stationary logic, or in second-order logic;
  \item splitting the relation variables \(\bar r\) as \(\bar s\bar t\) and
        freezing \(\bar s\) as parameters, which leaves pure logic and makes
        automorphisms, congruences and similar algebraic notions available;
  \item allowing the quantified relations to be families of subsets rather
        than subsets of some \(A^n\), giving definable third-order
        quantifiers; at this level the logic \(L(aa)\) is included;
  \item dropping definability and letting any class \(K\) of subsets of
        \(\bigcup_nA^n\) closed under isomorphism be a quantifier, as Shelah
        does, and likewise finite sequences of quantifiers
        \(\bar K=\langle K_0,\dots,K_n\rangle\).
\end{enumerate}

\begin{hbdefinition}{def:p03-w03-p03-046-b-expressible}{Expressibility (Definition~6.2)}
A finite sequence of quantifiers \(\bar K\) is \emph{expressible in} the
quantifier \(K\) if for each \(R\in K\) there is a formula
\(\phi(\bar x,\bar r)\), with quantifiers over the \(K_i\), such that for
some \(R_0,\dots,R_n\), each in one of the \(K_i\), \(R(\bar x)\iff
\phi(\bar x,\bar R)\) \cite[p.~477, Definition~6.2]{Baldwin1985DSOQ}.
\(K\) and \(\bar K\) are \emph{bi-expressible} if each is expressible in the
other.
\end{hbdefinition}

\begin{proofboundary}{def:p03-w03-p03-046-b-shelah-results}{Shelah's results reported in Section~6}
Baldwin reports two results of Shelah without proof
\cite[pp.~476--477]{Baldwin1985DSOQ}.
\begin{enumerate}
  \item (Theorem~6.1, from a paper of Shelah cited by Baldwin as
        [198?d].) For every first-order theory \(T\), the Hanf number of
        \((T,Q_{\mathrm{mon}})\) is at most \(H_{\mathrm{II}}\) iff
        \(Q_{\mathrm{II}}\not\le(T,Q_{\mathrm{mon}})\).
  \item (Theorem~6.3, from Shelah's [1983c].) The question, addressed in
        Shelah's [1983a], is whether every quantifier \(K\) is
        bi-interpretable with a finite sequence \(\bar K\) of quantifiers
        each of which is an equivalence relation. (i) If \(V=L\), every
        \(K\) is bi-expressible with an equivalence relation. (ii) It is
        consistent that some \(K\) is not bi-expressible with an
        equivalence relation.
\end{enumerate}
The proofs are in Shelah's papers, which the chapter cites but does not
reproduce; they are not given here.
\end{proofboundary}

Theorem~6.1 is the converse direction to the Hanf-number method of
Section~4: where that method uses a large Hanf number to show that
second-order logic is not interpreted, Shelah's theorem says that, for
monadic logic over a first-order theory, the absence of such an
interpretation already bounds the Hanf number by that of second-order
logic. Baldwin closes by noting that in the proofs of Vaught's conjecture
for \(\omega\)-stable theories and of Morley's conjecture on increasing
spectra, the step producing many models can be seen as an interpretation
of \(Q_{\mathrm{II}}\) in the \(L_{\omega_1\omega}(Q_{\mathrm{mon}})\)
theory of \(T\) \cite[p.~477]{Baldwin1985DSOQ}; that observation is
reported, not proved, here.

\section{Worked example and limitations}
\label{sec:p03-w03-p03-046-example}

\begin{hbexample}{def:p03-w03-p03-046-example}{Coding a relation by one involution}
Let \(T\) be the theory of an equivalence relation \(E\) with infinitely
many classes, all infinite, and let \(M=\omega\times\omega\) with
\((i,j)\mathrel E(i',j')\) iff \(i=i'\). The hypothesis of
\cref{thm:p03-w03-p03-046-equiv} holds with \(\phi=E\) and \(A=M\), so
\((T_\infty,\mathrm{2nd})\le(T,\text{1-1})\). The construction of the annex
proof, at \(\kappa=\aleph_0\), runs as follows. Take the classes
\(X_i=\{2i\}\times\omega\) for \(U_0\) and \(Y_j=\{2j+1\}\times\omega\) for
\(U_1\), and write \(x_{i,j}=(2i,j)\), \(y_{j,i}=(2j+1,i)\).
\begin{enumerate}
  \item The involution \(\alpha\) exchanges \(x_{i,j}\) and \(y_{j,i}\).
        Then the class \(X_i\) meets \(\alpha^{-1}[Y_j]\) in exactly the point
        \(x_{i,j}\).
  \item The domain of the interpreted set is \(B=\{x_{i,0}:i<\omega\}\),
        one point of each \(X_i\); the involution \(\varepsilon\) exchanges
        \(x_{i,0}\) and \(y_{i,0}\), and \(\gamma,\delta\) are involutions
        with \(\operatorname{Mov}(\gamma)=U_0\) and
        \(\operatorname{Mov}(\delta)=B\).
  \item A pair \((x_{i,0},x_{j,0})\in B\times B\) is represented by the
        unique \(x\in X_i\) with \(\alpha(x)\mathrel E\varepsilon(x_{j,0})\),
        namely \(x=x_{i,j}\).
  \item The relation \(R=\{(x_{0,0},x_{1,0})\}\) is coded by the involution
        \(\beta\) that exchanges \(x_{0,1}\) and \(y_{1,0}\) and fixes every
        other point, since \(\operatorname{Mov}(\beta)\cap U_0=\{x_{0,1}\}\).
        The successor relation \(\{(x_{i,0},x_{i+1,0}):i<\omega\}\) is coded
        by the involution exchanging each \(x_{i,i+1}\) with \(y_{i+1,i}\);
        it is the graph of an injective map of \(B\) into \(B\) that misses
        \(x_{0,0}\), and serves as the witness \(\zeta\) that \(B\) is
        infinite.
\end{enumerate}
\end{hbexample}

\begin{hbwarning}{def:p03-w03-p03-046-limits}{What these criteria do not give}
\begin{enumerate}
  \item The interpretations require clause~(2) of
        \cref{def:p03-w03-p03-046-interpretation} at every infinite
        cardinal. A first-order condition on the parameters cannot say that
        the interpreted domain is infinite, so each proof adds a coded
        injection that is not surjective; without it clause~(1) would fail
        at parameters defining a finite set.
  \item The lemmas give sufficient conditions only. They do not
        decide whether \((T_\infty,\mathrm{2nd})\le(T,\mathit{Mon})\) for a given
        theory without a coding, and the strictness assertions of the
        four-quantifier classification are not proved here.
  \item The monadic criteria quantify over arbitrary subsets of a model; the
        permutational criteria quantify over involutions, and a subset is
        available as a parameter only as the set of points moved, or fixed,
        by an involution.
\end{enumerate}
\end{hbwarning}

\emph{Relationships.} Generalized quantifiers of first-order type are
treated with partially ordered prefixes and the Henkin quantifier in
\cref{sec:p03-w03-p03-044-po-prefixes}. Existential second-order
sentences, which quantify over relations without a fixed quantifier
condition, appear in \cref{def:p03-existential-second-order}; that
definition concerns satisfaction of sentences in finite and infinite
structures, not the comparison of theories made here.
